% Pourquoi un leader ne valide pas une entrée d'un terme précédent par simple
% comptage (d'après la figure 8 d'Ongaro et Ousterhout, 2014).
\documentclass[border=6pt]{standalone}
\usepackage{sd}
\begin{document}
\begin{tikzpicture}[x=1mm, y=1mm, note/.append style={font=\sffamily\small},
    t1/.style={draw=Dslate, fill=DslateBg, line width=0.5pt, minimum size=6mm, font=\sffamily\scriptsize, inner sep=0pt},
    t2/.style={draw=Dblue, fill=DblueBg, line width=0.5pt, minimum size=6mm, font=\sffamily\scriptsize, inner sep=0pt},
    t3/.style={draw=Dorange, fill=DorangeBg, line width=0.5pt, minimum size=6mm, font=\sffamily\scriptsize, inner sep=0pt},
    t4/.style={draw=Dgreen, fill=DgreenBg, line width=0.5pt, minimum size=6mm, font=\sffamily\scriptsize, inner sep=0pt}]
  \newcommand{\jrn}[3]{% x, ligne (1..5), termes
    \foreach \t [count=\k] in {#3} \node[t\t] at ({#1+\k*7},{40-#2*8}) {\t};}
  \foreach \s in {1,...,5} \node[axe, anchor=east] at (2,{40-\s*8}) {S\s};
  % (a)
  \node[note, font=\sffamily\small\bfseries] at (14,40) {(a)};
  \jrn{0}{1}{1,2} \jrn{0}{2}{1,2} \jrn{0}{3}{1} \jrn{0}{4}{1} \jrn{0}{5}{1}
  \node[note, anchor=north, align=center, text width=42mm] at (14,-6) {S1, leader du terme 2,\\copie (2, terme 2) sur S2};
  % (b)
  \begin{scope}[xshift=46mm]
    \node[note, font=\sffamily\small\bfseries] at (14,40) {(b)};
    \jrn{0}{1}{1,2} \jrn{0}{2}{1,2} \jrn{0}{3}{1} \jrn{0}{4}{1} \jrn{0}{5}{1,3}
    \node[note, text=Dred] at (24,32) {×};
    \node[note, anchor=north, align=center, text width=42mm] at (14,-6) {S1 tombe ; S5 élu au terme 3\\(S3, S4, S5), écrit (2, terme 3)};
  \end{scope}
  % (c)
  \begin{scope}[xshift=92mm]
    \node[note, font=\sffamily\small\bfseries] at (14,40) {(c)};
    \jrn{0}{1}{1,2,4} \jrn{0}{2}{1,2} \jrn{0}{3}{1,2} \jrn{0}{4}{1} \jrn{0}{5}{1,3}
    \node[note, text=Dred] at (24,0) {×};
    \node[note, anchor=north, align=center, text width=42mm] at (14,-6) {S5 tombe ; S1 élu au terme 4,\\copie (2, terme 2) sur S3 :\\majorité, mais terme ancien};
  \end{scope}
  % (d)
  \begin{scope}[xshift=138mm]
    \node[note, font=\sffamily\small\bfseries] at (14,40) {(d)};
    \jrn{0}{1}{1,3} \jrn{0}{2}{1,3} \jrn{0}{3}{1,3} \jrn{0}{4}{1,3} \jrn{0}{5}{1,3}
    \node[note, anchor=north, align=center, text=Dred, text width=42mm] at (14,-6) {S1 tombe ; S5 élu (son dernier\\terme, 3, dépasse 2) et écrase\\la position 2 partout};
  \end{scope}
  \node[note, anchor=north, align=center] at (83,-26) {si S1 avait déclaré la position 2 validée en (c), une entrée validée aurait été effacée en (d) ;\\Raft ne valide par comptage que les entrées du terme courant (la position 3, terme 4, en (c))};
\end{tikzpicture}
\end{document}
